\section{Decision making under an uncertain physical state}\label{decision-making-under-an-uncertain-physical-state}

\subsection{1. Joint state and parameter uncertainty}\label{joint-state-and-parameter-uncertainty}

The physical forest map and the model coefficients are both uncertain. Plot,
remote-sensing, and operational records provide partial observations of the
physical state.

Let the latent physical state be

\[
X_t=(\mu_t,\nu_t,\kappa_t,r_t,\zeta_t),
\]

where the first three components describe tree, processor, and operator states,
\texttt{r\_t} is a discrete observation/coverage regime, and \texttt{zeta\_t} contains finite-
dimensional auxiliary states.

Let \texttt{theta} contain unknown structural parameters for growth, recovery,
processor acceptance, cost, and observations.

The controller observes a history of measurements and past controls. The joint
posterior is a sufficient state:

\[
\pi_t(dx,d\theta)
=P((X_t,\theta)\in dx\times d\theta\mid\mathcal I_t).
\]

This posterior carries uncertainty about both \textbf{what the system is} and \textbf{how
it evolves}.

\subsection{2. State and observation kernels}\label{state-and-observation-kernels}

Given action \texttt{u\_t}, the paper writes

\[
X_{t+1}\sim P_\theta(\cdot\mid X_t,u_t),
\]

\[
Y_{t+1}\sim G_\theta(\cdot\mid X_t,u_t,X_{t+1}).
\]

This separation is conceptually important. The transition kernel is about the
latent system. The observation kernel is about what is measured.

Examples of observations include:

\begin{itemize}
\tightlist
\item
  field plots;
\item
  processor intake and grade outcomes;
\item
  contractor productivity;
\item
  delivered-cost components;
\item
  repeated bids or procurement decisions;
\item
  remote-sensing features.
\end{itemize}

A satellite-derived index, for example, should enter \texttt{G\_theta} unless an
empirical observation model has established how it maps into the latent tree
state.

\subsection{3. Belief-state sufficiency}\label{belief-state-sufficiency}

The posterior predictive law of the next state and observation is obtained by
integrating the kernels against the current posterior. After observing
\texttt{Y\_\{t+1\}}, Bayes' rule produces \texttt{pi\_\{t+1\}}.

Thus the future conditional distribution depends on past history only through
\texttt{(pi\_t,u\_t)}. Under the stated Borel/measurability assumptions, any admissible
history-dependent policy has an equivalent policy measurable with respect to
the posterior.

This is the standard POMDP reduction in a very large state space.

For the conditional-law component, viscosity formulations on Wasserstein space
are available for jump diffusions, torus models, and fully second-order HJB
equations \citep{burzoni2019viscosity,soner2022torus,bayraktar2025hjb}.

\subsection{4. Bellman recursion on belief space}\label{bellman-recursion-on-belief-space}

Let \texttt{L(x,theta,u)} be one-step loss and

\[
\bar L(\pi,u)=\int L(x,\theta,u)\pi(dx,d\theta).
\]

If \texttt{K(dpi\textquotesingle{}\ \textbar{}\ pi,u)} is the Bayesian filtering kernel, the finite-horizon
recursion is

\[
V_t(\pi)=\inf_{u\in U(\pi)}
\left\{\bar L(\pi,u)+\gamma\int V_{t+1}(\pi')K(d\pi'\mid\pi,u)\right\}.
\]

Compact actions, lower-semicontinuous loss, weakly continuous belief
transitions, and a lower-semicontinuous continuation value give attainment of
the Bellman minimum.

The result establishes well-posedness under the stated conditions. Numerical
methods must still address the dimension of the belief state.

\subsection{5. The Wasserstein HJB subproblem}\label{the-wasserstein-hjb-subproblem}

For a smoother biological control subproblem, suppose the state is a probability
law \texttt{mu} on continuous tree coordinates and \texttt{V\_G(t,mu)} is differentiable in the
Lions sense.

The HJB contains:

\begin{itemize}
\tightlist
\item
  a time derivative;
\item
  a first measure derivative paired with drift;
\item
  a derivative in the state coordinate paired with local diffusion;
\item
  a second measure derivative generated by common noise.
\end{itemize}

Schematically,

\[
0=\partial_tV_G+
\inf_{a^G}\{F_G+\langle\partial_\mu V_G,b_G\rangle
+\text{diffusion term}+\text{common-noise second-order term}\}.
\]

The purpose of the equation in the paper is to expose the correct infinite-
dimensional derivative structure. The full belief-state market problem also
contains finite-mass geometry, hard thresholds, and impulses, so a single
smooth global HJB would be misleading.

\subsection{6. Nested dynamic risk}\label{nested-dynamic-risk}

Expected value alone can underweight severe downside. The paper uses a
one-step conditional risk map combining expectation and CVaR:

\[
\rho(Z)=(1-\lambda_R)E[Z]+\lambda_R\operatorname{CVaR}_{\alpha_R}(Z).
\]

Inserted recursively, this creates a time-consistent nested risk criterion.

For one-period loss,

\[
\operatorname{CVaR}_\alpha(Z)
=\min_z\left[z+\frac{1}{1-\alpha}E[(Z-z)_+]\right].
\]

This epigraph is convex when the underlying loss is convex in the decision.

\subsection{7. Distributional robustness is a separate layer}\label{distributional-robustness-is-a-separate-layer}

Bayesian posterior uncertainty is conditional on the model family. A separate
robustness layer addresses local structural misspecification.

The paper therefore defines an ambiguity set around the posterior predictive law
using Wasserstein distance:

\[
\mathcal P_{DR}(\pi,u)=\{Q:W_2(Q,\widehat P(\cdot\mid\pi,u))\le\epsilon_{DR}\}.
\]

A robust Bellman step replaces expected continuation value by a worst-case
expectation over that set.

\citet{esfahani2018dro} gives tractable reformulations and finite-sample
guarantees for Wasserstein ambiguity sets under standard regularity conditions.

Conceptually:

\begin{Shaded}
\begin{Highlighting}[]
\NormalTok{Bayesian uncertainty: which state/parameter inside the model is plausible?}
\NormalTok{Robust uncertainty:   what if the predictive model itself is locally wrong?}
\end{Highlighting}
\end{Shaded}

\subsection{8. Dual control and learning value}\label{dual-control-and-learning-value}

The action changes immediate economics and the distribution of future
observations. This creates an exploration-exploitation tradeoff.

Examples:

\begin{itemize}
\tightlist
\item
  inventorying a site to reduce uncertainty in harvestable size classes;
\item
  measuring processor intake to learn grade acceptance;
\item
  instrumenting contractor operations to identify productivity/cost;
\item
  offering a controlled contract to learn price response.
\end{itemize}

The paper defines an endogenous learning value by comparing continuation value
when the next observation is used versus ignored.

An information-theoretic experiment criterion such as expected information gain
can choose an experiment that learns parameters strongly. A decision-focused
criterion such as expected value of sample information can choose a different
experiment, because only uncertainty that changes the best action has direct
decision value.

Modern computational methods for Bayesian experimental design are reviewed by
\citet{rainforth2024bed}; variational estimators are developed by
\citet{foster2019boed}. The exploration effect is the dual-control mechanism
described by \citet{meijer2026dual}.

\subsection{9. A simple EVSI example}\label{a-simple-evsi-example}

Suppose there are two possible processor acceptance rates, low and high, and two
possible actions: commit to a large harvesting programme or stay small. An
experiment that tells the two states apart can be valuable only if the preferred
action differs between them.

If both parameter states imply the same best action, information gain can be
positive while EVSI is zero. The operational criterion therefore values the
effect of learning on the optimal action.

\subsection{10. Regime switching as impulse control}\label{regime-switching-as-impulse-control}

Let \texttt{R=\{r\_1,r\_2,r\_3\}} represent ordered observation/coverage regimes. Moving from
one regime to another incurs impulse cost \texttt{K\_\{r,r+\}} and maps the current
posterior into the enlarged regime through an embedding/prediction operator.

The quasi-variational recursion compares three choices:

\begin{enumerate}
\def\labelenumi{\arabic{enumi}.}
\tightlist
\item
  continue within the current regime;
\item
  pay the impulse cost and switch;
\item
  stop/abandon.
\end{enumerate}

The controller switches only when the value of the enlarged feasible
state/action region exceeds the value of waiting by enough to justify the cost.

This is the mathematical form of a staged evidence strategy: collect enough
information before committing to an expensive expansion.

\subsection{11. Computational reality}\label{computational-reality}

The exact posterior lives on a product of measure spaces and parameters. A
production controller will almost certainly need approximation: posterior
particles, parametric filters, low-dimensional embeddings, rollout, scenario
trees, or model-predictive control.

The structural theory is still useful because it tells us which summaries have
decision meaning. Candidate reduced belief coordinates include:

\begin{itemize}
\tightlist
\item
  local compatible supply by grade;
\item
  matched-demand fractions;
\item
  posterior moments of growth/quality;
\item
  dual matching values;
\item
  uncertainty around high-value specification thresholds.
\end{itemize}

The next chapters formalize the reduced-order projection and the data required
to identify these objects.
